\documentclass[11pt,a4paper]{article}

\usepackage[T1]{fontenc}
\usepackage[utf8]{inputenc}
\usepackage[turkish]{babel}
\usepackage{lmodern}
\usepackage{amsmath,amssymb,amsfonts}
\usepackage{mathtools}
\usepackage{bm}
\usepackage{siunitx}
\usepackage{geometry}
\usepackage{hyperref}
\usepackage{enumitem}
\usepackage{booktabs}
\usepackage{xcolor}

\geometry{margin=2.3cm}
\hypersetup{
  colorlinks=true,
  linkcolor=blue!50!black,
  citecolor=blue!50!black,
  urlcolor=blue!50!black,
  pdftitle={Prototipte Kullanılan BPG Algoritmasının Matematiksel Açıklaması},
  pdfauthor={OpenAI Codex}
}

\title{Prototipte Kullanılan BPG Algoritmasının\\Matematiksel Açıklaması}
\author{}
\date{\today}

\begin{document}
\maketitle

\begin{abstract}
Bu doküman, uygulamada \texttt{BPG} etiketiyle sunulan güdüm örneğinin gerçekte hangi matematiksel yapıyı kullandığını açıklar. En önemli nokta şudur: bu uygulamadaki \texttt{BPG}, literatürde tek ve standart bir kapalı form ile tanımlanan kanonik bir yasa değildir. Bunun yerine, \emph{pursuit}, \emph{body-pursuit} ve \emph{boresight-pursuit} ailesinden esinlenen, \texttt{gamma\_demand} üreten, mühendislik amaçlı sadeleştirilmiş bir prototiptir. Bu nedenle metin hem uygulamadaki gerçek denklemleri verir, hem de bunların pursuit, deviated pure pursuit ve proportional navigation ile ilişkisini açıklar.
\end{abstract}

\section{Özet Sonuç}

Uygulamadaki BPG örneğinin çekirdek formülü
\begin{equation}
\gamma_{\mathrm{cmd,raw}}
=
\gamma_m + k_1 \sigma + k_2 \dot{\lambda}
\label{eq:bpg_raw}
\end{equation}
şeklindedir. Burada
\begin{equation}
\sigma = \operatorname{wrap}_{[-\pi,\pi]}\!\left(\lambda - \gamma_m\right)
\label{eq:sigma_def}
\end{equation}
LOS açısı ile füze uçuş-yolu açısı arasındaki hatadır.

Dolayısıyla kullanılan yasa, en doğru ifadeyle,
\begin{quote}
\textbf{LOS yönelme hatasını ve LOS dönüş hızını kullanan, pursuit tabanlı, sezgisel bir flight-path shaping güdüm yasasıdır.}
\end{quote}

\section{Kodda Tam Olarak Ne Kullanılıyor?}

Uygulamadaki \texttt{BPG} örneği şu ifadeyle tanımlıdır:
\begin{equation}
\gamma_{\mathrm{cmd,raw}} = \gamma_m + k_1 \sigma + k_2 \dot{\lambda}.
\end{equation}

Örnek parametreler ilk açılış senaryosunda yaklaşık olarak
\begin{equation}
k_1 = 1.9, \qquad k_2 = 1.2, \qquad \tau_\gamma = 0.55 \ \si{s}
\end{equation}
olarak verilir.

Bu ham komut daha sonra uygulamada doğrudan kullanılmaz; önce güvenli bir açı aralığına kırpılır:
\begin{equation}
\gamma_{\mathrm{cmd}}
=
\operatorname{sat}\!\left(
\gamma_{\mathrm{cmd,raw}},
-
\left(\frac{\pi}{2}-0.05\right),
\frac{\pi}{2}-0.05
\right).
\label{eq:gamma_sat}
\end{equation}

Ardından birinci dereceden iç çevrim ile istenen açı hızına dönüştürülür:
\begin{equation}
\dot{\gamma}_{\mathrm{des}}
=
\frac{\gamma_{\mathrm{cmd}} - \gamma_m}{\tau_\gamma}.
\label{eq:gammadot_des}
\end{equation}

Gerekli normal ivme:
\begin{equation}
a_{z,\mathrm{need}}
=
V_m \dot{\gamma}_{\mathrm{des}} + g\cos\gamma_m
\label{eq:az_need}
\end{equation}
olarak hesaplanır ve füzenin ivme sınırı ile doyurulur:
\begin{equation}
a_z = \operatorname{sat}\!\left(a_{z,\mathrm{need}},-a_{\max},a_{\max}\right).
\label{eq:az_sat}
\end{equation}

Son olarak füze uçuş-yolu açısı dinamiği
\begin{equation}
\dot{\gamma}_m
=
\frac{a_z}{V_m} - \frac{g\cos\gamma_m}{V_m}
\label{eq:gamma_dyn}
\end{equation}
ile güncellenir.

Bu zincir çok önemlidir: uygulamadaki BPG doğrudan normal ivme komutu üreten bir yasa değil, önce \emph{uçuş-yolu açısı komutu} üreten, sonra bunu normal ivmeye çeviren bir yapıdır.

\section{2B Angajman Geometrisi}

Simülasyon \(x\)-\(z\) düzleminde kuruludur. Füze ve hedef konumları
\begin{align}
\bm{r}_m &= \begin{bmatrix}x_m \\ z_m\end{bmatrix},
&
\bm{r}_t &= \begin{bmatrix}x_t \\ z_t\end{bmatrix}
\end{align}
olsun.

Bağıl konum:
\begin{equation}
\bm{r} = \bm{r}_t - \bm{r}_m
=
\begin{bmatrix}
x_r \\ z_r
\end{bmatrix}
=
\begin{bmatrix}
x_t - x_m \\ z_t - z_m
\end{bmatrix}.
\end{equation}

Menzil:
\begin{equation}
R = \|\bm{r}\| = \sqrt{x_r^2 + z_r^2}.
\label{eq:R}
\end{equation}

LOS açısı:
\begin{equation}
\lambda = \operatorname{atan2}(z_r, x_r).
\label{eq:lambda}
\end{equation}

Füze ve hedef hız vektörleri:
\begin{align}
\bm{v}_m &= \begin{bmatrix}V_m\cos\gamma_m \\ V_m\sin\gamma_m\end{bmatrix},
&
\bm{v}_t &= \begin{bmatrix}V_t\cos\gamma_t \\ V_t\sin\gamma_t\end{bmatrix}.
\end{align}

Bağıl hız:
\begin{equation}
\bm{v}_r = \bm{v}_t - \bm{v}_m.
\end{equation}

Menzil türevi:
\begin{equation}
\dot{R}
=
\frac{x_r(v_{x,t}-v_{x,m}) + z_r(v_{z,t}-v_{z,m})}{R}.
\label{eq:Rdot}
\end{equation}

LOS dönüş hızı:
\begin{equation}
\dot{\lambda}
=
\frac{x_r(v_{z,t}-v_{z,m}) - z_r(v_{x,t}-v_{x,m})}{R^2}.
\label{eq:lambdadot}
\end{equation}

Bu ifade, 2B bağıl kinematiğin klasik sonucudur ve güdüm literatüründe temel büyüklüklerden biridir \cite{palumbo2010principles,li1999aopn}.

\section{Hata Değişkeni Olarak \texorpdfstring{\(\sigma\)}{sigma}}

Uygulamada
\begin{equation}
\sigma = \operatorname{wrap}_{[-\pi,\pi]}\!\left(\lambda-\gamma_m\right)
\end{equation}
kullanılır. Buradaki \(\operatorname{wrap}_{[-\pi,\pi]}\) işlemi, açıyı \([-\pi,\pi]\) aralığına sarar. Böylece:
\begin{itemize}[leftmargin=1.2cm]
  \item \(\sigma > 0\) ise LOS, füzenin mevcut uçuş doğrultusunun bir tarafındadır,
  \item \(\sigma < 0\) ise diğer tarafındadır,
  \item \(\sigma = 0\) ise füzenin hız vektörü LOS doğrultusuna hizalanmıştır.
\end{itemize}

Bu değişken geometrik olarak pursuit ailesinin doğal hata ölçüsüdür.

\section{BPG Yasasının Geometrik Yorumu}

Ham komut
\begin{equation}
\gamma_{\mathrm{cmd,raw}}
=
\gamma_m + k_1(\lambda-\gamma_m) + k_2\dot{\lambda}
\end{equation}
ifadesi açılırsa
\begin{equation}
\gamma_{\mathrm{cmd,raw}}
=
(1-k_1)\gamma_m + k_1 \lambda + k_2\dot{\lambda}
\label{eq:bpg_expanded}
\end{equation}
elde edilir.

Bu formüldeki terimler şu anlama gelir:
\begin{itemize}[leftmargin=1.2cm]
  \item \(k_1(\lambda-\gamma_m)\): füzenin yönünü LOS doğrultusuna çevirmeye çalışan pursuit geri beslemesi,
  \item \(k_2\dot{\lambda}\): LOS dönüyorsa buna erken tepki veren ek sönüm ya da faz-ileri terimi,
  \item \(\gamma_m\): komutu mevcut uçuş-yolu açısı etrafında inşa eden taşıyıcı terim.
\end{itemize}

Dolayısıyla bu yasa, saf pursuit ile LOS-rate katkısının bir karışımıdır.

\section{Pure Pursuit ile İlişkisi}

Eğer
\begin{equation}
k_1 = 1, \qquad k_2 = 0
\end{equation}
seçilirse
\begin{equation}
\gamma_{\mathrm{cmd,raw}} = \gamma_m + (\lambda-\gamma_m) = \lambda
\end{equation}
elde edilir. Yani özel durumda doğrudan \textbf{pure pursuit} komutuna inilir:
\begin{equation}
\gamma_{\mathrm{cmd}} = \lambda.
\end{equation}

Bu nedenle uygulamadaki BPG yasası, pure pursuit'in daha genel ve ayarlanabilir bir türevi gibi görülebilir.

\section{Deviated Pure Pursuit ile İlişkisi}

Deviated pure pursuit (DPP) literatüründe temel fikir, pursuer hız doğrultusunun LOS'a göre sıfır olmayan bir sapma ile tutulabilmesidir \cite{livermore2018dpp}. Bizim uygulamamızdaki yasa DPP'nin belirli bir optimal çözümünü uygulamaz; ancak aynı geometrik aileye aittir çünkü:
\begin{itemize}[leftmargin=1.2cm]
  \item kontrol değişkeni LOS ile uçuş doğrultusu arasındaki açı farkıdır,
  \item amaç doğrultu şekillendirmektir,
  \item komut doğrudan çarpışma geometri denklemi değil, yönelim temelli bir kuraldır.
\end{itemize}

Bu nedenle en doğru ifade şudur:
\begin{quote}
Bu prototipte kullanılan \texttt{BPG}, DPP ile akraba fakat onunla özdeş olmayan, body/boresight-pursuit esinli sezgisel bir gamma-komut yasasıdır.
\end{quote}

\section{İç Çevrim ve İvme Dinamiği}

Uygulama \texttt{gamma\_demand} modunda çalıştığı için dış güdüm çevrimi ile iç uçuş kontrol çevrimi ayrılmıştır.

\subsection{Normal ivme yorumu}

Denklem \eqref{eq:az_need}, sabit hız varsayımında klasik
\begin{equation}
a_n \approx V \dot{\gamma}
\end{equation}
ilişkisine, yerçekimi düzeltmesi eklenmiş hal olarak okunabilir. Füze dikey düzlemde uçtuğu için yerçekiminin uçuş-yolu açısı üzerindeki etkisi ayrıca hesaba katılır \cite{jackson2010fcs}.

\subsection{Doyumsuz bölgede yaklaşık kapalı çevrim}

Eğer \(a_z\) doyuma girmiyorsa, \eqref{eq:az_need} ve \eqref{eq:gamma_dyn} birlikte yaklaşık olarak
\begin{equation}
\dot{\gamma}_m
=
\frac{\gamma_{\mathrm{cmd}}-\gamma_m}{\tau_\gamma}
\end{equation}
verir.

BPG yerleştirilirse:
\begin{equation}
\dot{\gamma}_m
=
\frac{k_1\sigma + k_2\dot{\lambda}}{\tau_\gamma}.
\label{eq:gammadot_cl}
\end{equation}

Çünkü \(\sigma = \lambda - \gamma_m\), hata dinamiği yaklaşık olarak
\begin{align}
\dot{\sigma}
&=
\dot{\lambda} - \dot{\gamma}_m \\
&=
\dot{\lambda}
-
\frac{k_1\sigma + k_2\dot{\lambda}}{\tau_\gamma} \\
&=
-\frac{k_1}{\tau_\gamma}\sigma
+
\left(1-\frac{k_2}{\tau_\gamma}\right)\dot{\lambda}
\label{eq:sigma_dyn}
\end{align}
şeklinde yazılır.

Bu ifade bize şunu söyler:
\begin{itemize}[leftmargin=1.2cm]
  \item \(k_1 > 0\) oldukça \(\sigma\) hatası bastırılır,
  \item \(k_2\) LOS dönüş hızına doğrudan tepki vererek pure pursuit gecikmesini azaltır,
  \item \(\tau_\gamma\) küçüldükçe iç çevrim daha çevik davranır.
\end{itemize}

\section{Katsayıların Fiziksel Yorumu}

\subsection{\texorpdfstring{\(k_1\)}{k1}}

\(k_1\), LOS hizalama hatasının kazancıdır:
\begin{itemize}[leftmargin=1.2cm]
  \item \(k_1 = 1\): pure pursuit benzeri davranış,
  \item \(k_1 < 1\): LOS'a daha yumuşak yönelme,
  \item \(k_1 > 1\): LOS'a daha agresif kapanma.
\end{itemize}

\subsection{\texorpdfstring{\(k_2\)}{k2}}

\(\dot{\lambda}\) birimi \si{rad/s} olduğu için, boyutsal açıdan \(k_2\)'nin birimi saniye gibi yorumlanabilir. Uygulamada kullanıcı bunu skaler sayı olarak girer; fakat fiziksel etkisi bir zaman-ölçeği gibidir:
\begin{itemize}[leftmargin=1.2cm]
  \item \(k_2 = 0\): saf pursuit karakteri baskın,
  \item \(k_2 > 0\): dönen LOS'a daha erken tepki,
  \item çok büyük \(k_2\): gereksiz agresiflik ve gürültü hassasiyeti.
\end{itemize}

\subsection{\texorpdfstring{\(\tau_\gamma\)}{tau-gamma}}

\(\tau_\gamma\), iç çevrim takip zaman sabitidir:
\begin{itemize}[leftmargin=1.2cm]
  \item küçük \(\tau_\gamma\): hızlı takip,
  \item büyük \(\tau_\gamma\): daha yumuşak fakat daha gecikmeli takip.
\end{itemize}

\section{PNG ile Farkı}

Uygulamadaki PNG örneği
\begin{equation}
a_{z,\mathrm{PNG}} = N V_m \dot{\lambda}
\label{eq:png}
\end{equation}
şeklindedir.

PNG'nin ana mantığı LOS dönüş hızını sıfıra sürükleyerek çarpışma geometrisini korumaktır \cite{palumbo2010principles}. Buna karşılık uygulamadaki BPG,
\begin{equation}
\gamma_{\mathrm{cmd}} = \gamma_m + k_1\sigma + k_2\dot{\lambda}
\end{equation}
ile esas olarak doğrultu şekillendirme problemi çözer.

\begin{center}
\begin{tabular}{@{}lll@{}}
\toprule
Özellik & PNG & Prototip BPG \\
\midrule
Temel geri besleme & \(\dot{\lambda}\) & \(\sigma\) ve \(\dot{\lambda}\) \\
Komut tipi & normal ivme & uçuş-yolu açısı komutu \\
Karakter & lead/collision odaklı & pursuit/pointing odaklı \\
İç çevrim bağımlılığı & daha düşük & daha yüksek \\
\bottomrule
\end{tabular}
\end{center}

\section{Sınırlamalar ve Doğru Konumlandırma}

Bu dokümanda anlatılan yasa aşağıdaki nedenlerle sezgiseldir:
\begin{enumerate}[leftmargin=1.2cm]
  \item Tek bir standart yayınlanmış BPG formülünün aynısı değildir.
  \item Füze 2B point-mass modeldir; seeker, airframe ve gerçek otopilot ayrıntıları yoktur.
  \item Komut zincirinde hem açı kırpma hem normal ivme doyumu vardır.
  \item Hedef ilk sürümde sabit hız ve sabit \(\gamma_t\) ile uçmaktadır.
  \item \(k_1\), \(k_2\) ve \(\tau_\gamma\) değerleri senaryoya bağlıdır; evrensel optimum değildir.
\end{enumerate}

Bu yüzden uygulamadaki \texttt{BPG} en iyi şu şekilde adlandırılmalıdır:
\begin{quote}
\textbf{Pursuit tabanlı, body/boresight esinli, gamma-komut üreten prototip bir güdüm yasası.}
\end{quote}

\section{Sonuç}

Uygulamadaki BPG örneğinin matematiksel çekirdeği
\begin{equation}
\gamma_{\mathrm{cmd,raw}}
=
\gamma_m + k_1(\lambda-\gamma_m) + k_2\dot{\lambda}
\end{equation}
olup, bu komut daha sonra doyurulur, iç çevrim ile \(\dot{\gamma}\) isteğine çevrilir ve normal ivme sınırı ile uygulanır.

Dolayısıyla bu algoritma:
\begin{enumerate}[leftmargin=1.2cm]
  \item LOS-hizalama hatasını azaltır,
  \item LOS dönüş hızına tepki verir,
  \item doğrudan çarpışma geometrisi yerine yönelim şekillendirme yapar.
\end{enumerate}

Eğer \(k_1 = 1\) ve \(k_2 = 0\) seçilirse pure pursuit'e iner. Eğer LOS-rate katkısı eklenirse pure pursuit'in gecikmesi azaltılmaya çalışılır. Bu nedenle uygulamadaki BPG, eğitim ve hızlı deney için uygun, fakat literatürdeki tek bir kanonik yasa ile birebir özdeş olmayan bir mühendislik prototipidir.

\begin{thebibliography}{9}

\bibitem{watkins1946}
C. E. Watkins,
\newblock \emph{Paths of Target Seeking Missiles in Two Dimensions},
\newblock NACA ACR L6B06, 1946.
\newblock NASA NTRS kayıt sayfası:
\newblock \url{https://ntrs.nasa.gov/citations/20090033671}

\bibitem{palumbo2010principles}
N. F. Palumbo, R. A. Blauwkamp ve J. M. Lloyd,
\newblock ``Basic Principles of Homing Guidance,''
\newblock \emph{Johns Hopkins APL Technical Digest}, Cilt 29, Sayı 1, 2010, ss. 25--41.
\newblock PDF:
\newblock \url{https://www.jhuapl.edu/Content/techdigest/pdf/V29-N01/29-01-Palumbo_Principles_Rev2018.pdf}

\bibitem{jackson2010fcs}
P. B. Jackson,
\newblock ``Overview of Missile Flight Control Systems,''
\newblock \emph{Johns Hopkins APL Technical Digest}, Cilt 29, Sayı 1, 2010.
\newblock PDF:
\newblock \url{https://www.jhuapl.edu/sites/default/files/2024-09/29-01-Jackson.pdf}

\bibitem{li1999aopn}
M. Li ve C.-C. Lim,
\newblock ``Observability-Enhanced Proportional Navigation Guidance with Bearings-Only Measurements,''
\newblock \emph{ANZIAM Journal}, 1999.
\newblock PDF:
\newblock \url{https://www.cambridge.org/core/services/aop-cambridge-core/content/view/491048E9EF20684EFB5FB8F4703CD069/S0334270000010584a.pdf/observabilityenhanced_proportional_navigation_guidance_with_bearingsonly_measurements.pdf}

\bibitem{livermore2018dpp}
R. Livermore ve T. Shima,
\newblock ``Deviated Pure-Pursuit-Based Optimal Guidance Law for Imposing Intercept Time and Angle,''
\newblock \emph{Journal of Guidance, Control, and Dynamics}, Cilt 41, Sayı 8, 2018, ss. 1805--1811.
\newblock Açık erişim kopya:
\newblock \url{https://aerospace.technion.ac.il/lab/casy/wp-content/uploads/sites/12/2024/01/livermore-shima-2018-deviated-pure-pursuit-based-optimal-guidance-law-for-imposing-intercept-time-and-angle.pdf}

\end{thebibliography}

\end{document}
